\chapter{Classical Controller Implementation}
\label{app:classical_controllers}

The deep reinforcement-learning policy is benchmarked against four classical
path-tracking controllers. This appendix documents their implementation: the exact
control laws, the forward and reverse variants, the way each controller was tuned, and
the harness used to evaluate them. The four controllers are chosen to span the classical
control spectrum, from a purely geometric law through to constrained optimal control with
preview:

\begin{itemize}
    \item \textbf{Pure Pursuit} --- a geometric proportional law;
    \item \textbf{PID} --- classical error feedback with an integral term;
    \item \textbf{LQR} --- optimal static linear feedback on the error dynamics;
    \item \textbf{Kinematic LTV MPC} --- optimal constrained control with preview.
\end{itemize}

All four are implemented in the batched simulator
(\texttt{tractor\_trailer\_rl\_cupy}) so that they run inside the same vectorized
environment as the learned policies. Crucially, every controller emits the same action
as the RL agent --- a bounded \emph{steering rate} --- so the comparison between
classical and learned control is made on identical terms.

\section{Common framework}
\label{sec:app_cc_common}

\subsection{Observation layout and shared action mapping}

Each controller is a function of the observation vector already produced by the
environment. For the trailer tasks this vector is
\begin{equation}
    o = [\,s,\ \gamma,\ e_y,\ e_\psi,\ e_{y,t},\ e_{\psi,t},\ \kappa,\ \dots\,],
    \label{eq:app_cc_obs}
\end{equation}
where $s$ is the current steering angle, $\gamma$ the hitch (articulation) angle, and the
error pairs $(e_y, e_\psi)$ and $(e_{y,t}, e_{\psi,t})$ are the cross-track and heading
errors of the tractor and of the trailer respectively; $\kappa$ is the reference path
curvature at the lookahead point. The trailing entries carry additional task-specific
channels (for example lidar returns) that the classical controllers do not use.

Every controller computes a desired steering angle $\delta_{\mathrm{des}}$ and then
converts it to the environment's steering-rate action through a common saturating map,
\begin{equation}
    \dot{\delta} = \operatorname{clip}\!\left(\frac{\delta_{\mathrm{des}} - s}{\Delta t},\
    -\dot{\delta}_{\max},\ +\dot{\delta}_{\max}\right),
    \label{eq:app_cc_steerrate}
\end{equation}
where $\Delta t$ is the simulation step and $\dot{\delta}_{\max}$ the steering-rate limit.
Because the learned policy also outputs $\dot{\delta}$, and \cref{eq:app_cc_steerrate}
applies to all four baselines identically, the classical and learned controllers act
through exactly the same interface.

\subsection{Forward and reverse tracking}
\label{sec:app_cc_fwdrev}

Each controller has a forward and a reverse variant, reflecting a structural asymmetry of
the articulated vehicle. Driving \emph{forward}, the tractor leads and the trailer follows
passively, so it suffices to regulate the tractor errors $(e_y, e_\psi)$ towards zero.
Driving \emph{in reverse}, the trailer leads and the coupling inverts: the hitch angle
$\gamma$ becomes \emph{open-loop unstable} (small articulation errors grow and the vehicle
jackknifes), so the reverse controllers regulate the trailer errors
$(e_{y,t}, e_{\psi,t})$ and must additionally include an explicit hitch-stabilizing action.
This asymmetry is why the reverse variants below are not simply the forward laws with a
sign change, and why their gains had to be tuned separately
(\cref{sec:app_cc_tuning}).

\subsection{Vehicle parameters}

The benchmark uses the full-size shunt-truck configuration. The parameters that enter the
controllers are summarized in \cref{tab:app_cc_params}; here $L_1$ is the tractor
wheelbase, $L_2$ the effective trailer length (hitch to trailer axle), and $V$ the fixed
longitudinal speed held constant throughout the fixed-speed path-tracking experiments. The
large length ratio $L_2/L_1 \approx 4.2$ is what makes the reverse hitch dynamics
difficult to stabilize.

\begin{table}[H]
\centering
\caption{Shunt-truck parameters used by the classical controllers.}
\label{tab:app_cc_params}
\begin{tabular}{llc}
\toprule
Symbol & Quantity & Value \\
\midrule
$L_1$ & Tractor wheelbase & $2.95$\,m \\
$L_2$ & Trailer length (hitch to axle) & $12.5$\,m \\
$V$ & Fixed longitudinal speed & $5.0$\,m/s \\
$\Delta t$ & Simulation step & $0.1$\,s \\
$\delta_{\max}$ & Steering-angle limit & $0.87$\,rad ($\approx 50^{\circ}$) \\
$\dot{\delta}_{\max}$ & Steering-rate limit & $25^{\circ}$/s \\
\bottomrule
\end{tabular}
\end{table}

In the reverse variants the speed enters with a sign, $V < 0$. Because the environment's
tractor yaw response to steering does not itself flip sign in reverse, the along-path
kinematics use the magnitude $\lvert V \rvert$ while the yaw dynamics retain the signed
$V$; this distinction is carried consistently through the LQR and MPC models below.

\section{Pure Pursuit}
\label{sec:app_cc_pp}

The pure-pursuit controller \cite{coulter1992pure} is a fixed proportional geometric law.
Forward, it combines a curvature feed-forward with proportional cross-track and heading
feedback on the tractor errors,
\begin{equation}
    \delta_{\mathrm{des}} = -\bigl(k_{\mathrm{ff}} \arctan(L_1 \kappa)
    + k_y\, e_y + k_\theta\, e_\psi\bigr).
    \label{eq:app_cc_pp_fwd}
\end{equation}
Reverse, it regulates the trailer errors and adds an explicit hitch-stabilizing term
$k_{\mathrm{hitch}}\,\gamma$ to counter the open-loop-unstable articulation,
\begin{equation}
    \delta_{\mathrm{des}} = k_{\mathrm{hitch}}\,\gamma + k_{\mathrm{ff}}\,\kappa
    + k_y\, e_{y,t} + k_\theta\, e_{\psi,t}.
    \label{eq:app_cc_pp_rev}
\end{equation}
Both laws are stateless functions of the observation, so the controller is fully
vectorized across the parallel environments. Besides serving as a benchmark, the tuned
pure-pursuit controller doubles as the reference (``teacher'') signal for the guided
reward and as a feasibility oracle: a path that a well-tuned pure-pursuit controller can
complete is deemed geometrically feasible.

\section{PID}
\label{sec:app_cc_pid}

The PID controller extends the proportional law with an integral term that removes
steady-state cross-track offset. The integral is accumulated per environment and clamped
to prevent windup, and it is reset at episode boundaries. Forward,
\begin{equation}
    \delta_{\mathrm{des}} = -k_{\mathrm{ff}}\,\kappa - K_p\, e_y
    - K_i \!\int\! e_y - K_d\, e_\psi + K_{p,t}\, e_{y,t},
    \label{eq:app_cc_pid_fwd}
\end{equation}
and reverse (regulating the trailer errors, with the hitch stabiliser),
\begin{equation}
    \delta_{\mathrm{des}} = k_{\mathrm{hitch}}\,\gamma + k_{\mathrm{ff}}\,\kappa
    + K_p\, e_{y,t} + K_i \!\int\! e_{y,t} + K_d\, e_{\psi,t}.
    \label{eq:app_cc_pid_rev}
\end{equation}
The derivative gain $K_d$ multiplies the heading error, so it acts as a heading-damping
term rather than a numerical time derivative. The integrator uses the anti-windup update
\begin{equation}
    I \leftarrow \operatorname{clip}\!\bigl(I + e\,\Delta t,\ -I_{\max},\ +I_{\max}\bigr),
    \label{eq:app_cc_pid_int}
\end{equation}
where $e$ is the regulated cross-track error ($e_y$ forward, $e_{y,t}$ reverse). Because
the controller carries this per-environment state, the evaluation loop passes a
termination mask so that $I$ is zeroed for any environment that resets on the current
step.

\section{LQR}
\label{sec:app_cc_lqr}

The LQR baseline is the canonical optimal-linear controller: static feedback $u = -Kx$ on
a linear model of the tracking-error dynamics, plus a curvature feed-forward
\cite{rajamani2011vehicle}. The error state is
\begin{equation}
    x = [\,e_{\mathrm{lat}},\ e_{\mathrm{head}},\ \gamma\,]^\top,
    \label{eq:app_cc_lqr_state}
\end{equation}
the lateral and heading error of the leading unit (tractor forward, trailer reverse)
together with the hitch angle; the input is the steering angle $u = \delta$. The
continuous error dynamics $\dot{x} = A_c x + B_c u$ are direction-specific. Forward,
\begin{equation}
    A_c = \begin{bmatrix}
        0 & V & 0 \\[2.2mm]
        0 & 0 & 0 \\[2.2mm]
        0 & 0 & -\dfrac{V}{L_2}
    \end{bmatrix},
    \qquad
    B_c = \begin{bmatrix}
        0 \\[2.2mm] \dfrac{V}{L_1} \\[2.2mm] \dfrac{V}{L_1}
    \end{bmatrix},
    \label{eq:app_cc_lqr_fwd}
\end{equation}
and reverse, where the leading trailer progresses along the path at speed
$\lvert V \rvert$ while the yaw dynamics keep the signed $V$,
\begin{equation}
    A_c = \begin{bmatrix}
        0 & \lvert V \rvert & 0 \\[2.2mm]
        0 & 0 & \dfrac{V}{L_2} \\[2.2mm]
        0 & 0 & -\dfrac{V}{L_2}
    \end{bmatrix},
    \qquad
    B_c = \begin{bmatrix}
        0 \\[2.2mm] 0 \\[2.2mm] \dfrac{\lvert V \rvert}{L_1}
    \end{bmatrix}.
    \label{eq:app_cc_lqr_rev}
\end{equation}
The pair $(A_c, B_c)$ is discretized by zero-order hold at the step $\Delta t$ to give
$(A_d, B_d)$. With state and input weights $Q = \operatorname{diag}(q_{\mathrm{lat}},
q_{\mathrm{head}}, q_\gamma)$ and $R$, the infinite-horizon gain follows from the discrete
algebraic Riccati equation. Its stabilizing solution $P$ yields
\begin{equation}
    K = \bigl(R + B_d^\top P B_d\bigr)^{-1} B_d^\top P A_d.
    \label{eq:app_cc_lqr_gain}
\end{equation}
Because the speed is constant, $K$ is computed once at construction and then applied to
every environment, so the controller stays fully vectorized. The command adds a
steady-state feed-forward that holds the vehicle on a constant-curvature arc: a
feed-forward steer $\delta_{\mathrm{ff}} = L_1 \kappa$ and a hitch reference
$\gamma_{\mathrm{ref}} = \operatorname{sign}(V)\, L_2\, \kappa$, giving
\begin{equation}
    \delta = \delta_{\mathrm{ff}} - K\,(x - x_{\mathrm{ref}}),
    \qquad
    x_{\mathrm{ref}} = [\,0,\ 0,\ \gamma_{\mathrm{ref}}\,]^\top,
    \label{eq:app_cc_lqr_cmd}
\end{equation}
which is then saturated to $\pm\delta_{\max}$ and mapped to a steering rate through
\cref{eq:app_cc_steerrate}. The reverse weights place a heavy penalty on control effort
and hitch deviation, which is what damps the otherwise unstable reverse plant.

\section{Kinematic LTV MPC}
\label{sec:app_cc_mpc}

The final baseline is a linear time-varying model-predictive controller
\cite{rawlings2017mpc} that adds a finite preview horizon and explicit input constraints.
The MPC uses the same error state $x = [\,e_{\mathrm{lat}}, e_{\mathrm{head}}, \gamma\,]^\top$
as the LQR but plans over an $N$-step horizon.

\subsection{Prediction models}

Forward, the tractor leads, so the MPC tracks the tractor errors with the steering angle
$\delta$ as the input, giving the same $(A_c, B_c)$ as the forward LQR model
\eqref{eq:app_cc_lqr_fwd}. Reverse, servoing the trailer pose with the unstable hitch
dynamics inside the optimization causes the long shunt trailer to oscillate and
jackknife. The reverse controller therefore uses a flatness / virtual-input
reformulation (following the Autoware trailer LTV-MPC): the QP plans the \emph{trailer} as
a stable bicycle driven by a \emph{virtual} trailer steer $\delta_T$, and the hitch enters
only through an inner loop that has been closed to the stable first-order dynamics
$\dot{\gamma} = K(\delta_T - \gamma)$,
\begin{equation}
    A_c = \begin{bmatrix}
        0 & \lvert V \rvert & 0 \\[2.2mm]
        0 & 0 & \dfrac{V}{L_2} \\[2.2mm]
        0 & 0 & -K
    \end{bmatrix},
    \qquad
    B_c = \begin{bmatrix}
        0 \\[2.2mm] 0 \\[2.2mm] K
    \end{bmatrix}.
    \label{eq:app_cc_mpc_rev}
\end{equation}
The reference curvature enters both models as a measured disturbance through
$E_c = [\,0,\ -\lvert V \rvert,\ 0\,]^\top$. After the QP returns the virtual input, an
analytic inner map realises it on the physical tractor and supplies the anti-jackknife
counter-steer,
\begin{equation}
    \delta_f = \arctan\!\left(\frac{L_1}{L_2}\,\operatorname{sign}(V)\,\sin\gamma
    + \frac{L_1 K}{\lvert V \rvert}\,(\delta_T - \gamma)\right).
    \label{eq:app_cc_mpc_inner}
\end{equation}
The $1/\lvert V \rvert$ factor is what flips the feedback sign in reverse and produces the
stabilizing counter-steer.

\subsection{Condensed quadratic program}

The continuous models are discretized by zero-order hold and the state trajectory is
condensed onto the input sequence, so that the optimization variable is the vector of
$N$ future steering inputs. With the stacked weight $\bar{Q} = I_N \otimes Q$, a control
penalty $R$, an input-rate penalty $R_d$, and the first-difference operator
$T = I - I^{(-1)}$ (so that $T$ forms $\Delta u$), the cost Hessian is
\begin{equation}
    H = 2\bigl(B_u^\top \bar{Q}\, B_u + R\, I + R_d\, T^\top T\bigr),
    \label{eq:app_cc_mpc_cost}
\end{equation}
where $B_u$ is the condensed input-to-state map. The problem is subject to a steering-rate
box on the increments, $\lvert \Delta u \rvert \le \dot{\delta}_{\max}\,\Delta t$, and a
steering box $\lvert u \rvert \le \delta_{\max}$. It is solved per environment with OSQP.
The rate-penalty term $R_d\,T^\top T$ is weighted heavily in reverse so the controller
does not chase cross-track measurement ripple. The overall per-step procedure is
summarized in \cref{alg:app_cc_mpc}.

\begin{algorithm}[H]
\caption{Kinematic LTV MPC step (per environment)}
\label{alg:app_cc_mpc}
\begin{algorithmic}[1]
\State \textbf{Precompute (once):} discretize $(A_c, B_c, E_c)$ by ZOH; build the
condensed maps $A_x, B_u, E_w$ and the constant Hessian $H$ from \eqref{eq:app_cc_mpc_cost}.
\State \textbf{On episode reset:} clear the warm-start input $u_{\text{prev}} \gets 0$.
\State Read $x_0 = [\,e_{\mathrm{lat}}, e_{\mathrm{head}}, \gamma\,]$ and curvature $\kappa$
from the observation.
\State Form the linear cost term from $x_0$, the previewed curvature disturbance, and
$u_{\text{prev}}$.
\State Solve the QP with OSQP subject to the steering-rate and steering boxes.
\If{solver fails} $u \gets u_{\text{prev}}$ \Comment{safe fallback}
\Else{} $u \gets$ first element of the optimal input sequence
\EndIf
\State $u_{\text{prev}} \gets u$
\State $\delta \gets u$ (forward) \textbf{or} apply inner map \eqref{eq:app_cc_mpc_inner}
to obtain $\delta_f$ (reverse).
\State \Return steering rate from \eqref{eq:app_cc_steerrate}.
\end{algorithmic}
\end{algorithm}

Only the first input of the optimal sequence is applied (receding horizon); the solved
input is retained as the warm start for the next step and, like the PID integrator, is
reset at episode boundaries. If the solver fails to converge, the controller falls back to
the previous command.

\section{Gain-tuning methodology}
\label{sec:app_cc_tuning}

The controllers above are defined by their gains and weights: the proportional and
feed-forward gains of pure pursuit, the PID gains, the LQR state/input weights, and the
MPC horizon and cost weights. Rather than report a specific set of numbers, which are
particular to one vehicle configuration, the point of interest is the procedure used to
obtain them, which is the same for every controller.

Each controller is tuned by a bounded random search (a gain sweep) run directly against
the simulator \cite{storn1997differential}. On each iteration a candidate gain set is
sampled uniformly from hand-set ranges, the controller is rolled out over a batch of
parallel environments, and the candidate is scored by the lexicographic key
\begin{equation}
    \bigl(\text{completion rate},\ -\,\text{mean cross-track error}\bigr),
    \label{eq:app_cc_tuning_key}
\end{equation}
so that reliably completing the path is maximized first and, among gains that complete,
the one with the tightest tracking is preferred. The best candidate is retained. This
search is run \emph{separately for each controller and each direction}, because the
forward and reverse plants differ so strongly.

The reverse variants in particular could not simply inherit the gains tuned on the small
laboratory vehicle: on the shunt truck those gains destabilize the vehicle, and several
gains (notably the hitch-stabilizing terms) change sign. This is a direct consequence of
the open-loop-unstable reverse hitch dynamics and the large trailer-to-tractor length
ratio $L_2/L_1 \approx 4.2$ discussed in \cref{sec:app_cc_fwdrev}; the reverse controllers
were therefore re-tuned from scratch for the target geometry.

\section{Evaluation harness}
\label{sec:app_cc_eval}

All four controllers are evaluated by a single harness that reuses the same batched
environment as the reinforcement-learning policies, so that the classical baselines and
the learned policy are measured on identical paths under identical dynamics. The harness
runs both driving directions and, for each controller, records three per-episode metrics:

\begin{itemize}
    \item \textbf{Completion} --- the fraction of episodes that reach the goal
    (the mean success flag);
    \item \textbf{Trailer cross-track error} --- the mean absolute lateral error of the
    trailer over the episode, a measure of tracking accuracy;
    \item \textbf{Peak hitch angle} --- the largest articulation angle reached, a measure
    of how close the vehicle came to jackknifing (the jackknife threshold is
    $1.4$\,rad).
\end{itemize}

The harness handles the stateful controllers by passing a per-environment termination mask
into the controller at each step, so the PID integrator \eqref{eq:app_cc_pid_int} and the
MPC warm start are reset exactly when an environment restarts. Evaluating every controller
in the same loop, on the same feasible path pool, is what makes the reinforcement-learning
comparison in the results chapter a fair, like-for-like one. The measured numbers are
reported there rather than in this appendix.
